\documentclass{article}
\usepackage{RNotation}

\title{Quantum mechanics}
\author{Ross Grogan}
\date{}

\begin{document}
	
\maketitle

\section*{The double-slit experiment}

Huygens and Fresnel discovered that if we perform a double slit experiment with water waves, a single wave emitted by a source will be split into two copies of itself by the barrier between the slits, so that each slit becomes an emitter. We can repeat this experiment with light, and see interference patterns perfectly predicted by the classical model.

\begin{itemize}
	\item (Strange new observation). Except, by turning the intensity of the source down, we observe that light is detected on the screen as discrete particles- nothing like a wave hitting a breakwater. And if we have a single-photon source, we can further observe that exactly one photon is detected at a time. (If we don’t have a single photon source, we can still deduce that one photon is detected at a time by placing a photon detector at each slit, and seeing that the two slits never both record a photon.)
\end{itemize}

So, we have:

\begin{axiom}
	(Wave interference). Light functionally travels like a wave that comes out of the emitter as one wave and is split by the “breakwater” into two identical waves that interfere with each other and sum to form an aggregate wave.
\end{axiom}

\begin{axiom}
	(Discreteness). The aggregate wave arrives as exactly one of the particles that would have constituted it.
\end{axiom}

And, since we observe the same interference pattern as in the classical case:

\begin{axiom}
	(Axiom: structured randomness). The location of that particle is randomly chosen, such that the distribution of the probabilities of the possible locations is proportional to the intensity of the aggregate wave.
\end{axiom}

One more observation yields some more unexpected behavior. If we place photon detectors at each slit, the interference patter disappears, and we only see brightness. To account for this, we propose another axiom:

\begin{axiom}
	(Indistinguishable alternatives). Only reemitted waves corresponding to indistinguishable alternative paths are included in the aggregate wave.
\end{axiom}

Let me briefly explain how the observation follows from this axiom. [TO-DO]. While there are potentially other axioms that could explain this most recent observation, this one seems relatively natural.

\section*{Pedagogy}

\begin{itemize}
	\item I know that somehow, one is supposed to be able to conclude wave-particle duality; that one can think of light as a particle with wave characteristics just as much as one can think of it as a wave with particle characteristics. But, all we can be sure of from this experiment is that a particle is detected on the screen. How can we be sure that it was a particle when it was emitted, and when it traveled? We have no way of detecting that. I suspect more experiments are necessary for uncovering the full wave-particle duality.
	\item At least for this experiment it seems that thinking of light as a wave seems to explain things better than thinking of it as a particle does. It's instinctive to imagine particles taking definite paths through space, and therefore easy to think that photons have merely paths that vary wildly as initial conditions change, but are nevertheless known by the universe. It's even easy to hold onto this idea of definitness while still permitting the strange new possibility for photons to experience interference. Before one fires one photon at a time, the situation could be "interference, but no randomness"- maybe photons interfere with each other rather than with themselves. And even when one does fire one photon at a time, it's very difficult to come up with explanations for what happened to the photon in-between the firing and the landing on the screen. In fact, how can we even be sure if the light was a photon when it was fired?
\end{itemize}


\end{document}